\section{Analysis}

\subsection{Encoder vs Decoder Dispatch}

\textbf{Observation.} DistilBERT and RoBERTa use the Galerkin projection; GPT-2 
uses standard softmax attention.

\textbf{Why?} Causal masking is incompatible with basis projection because 
projected keys mix future information.

\textbf{Implication.} Galerkin provides no benefit on GPT-2.

\subsection{Basis Approximation Quality}

\textbf{Figure 8.1} shows the spectrum of DCT coefficients for a sample attention 
output. The top 64 coefficients capture ~92\% of the energy, 
justifying $r = 64$.

\subsection{Per-Layer Reconstruction Error}

\textbf{Table 8.1: Frobenius error of attention output vs full softmax (DistilBERT).}

\begin{table}[H]
\centering
\small
\begin{tabular}{ll}
\toprule
\textbf{Layer} & \textbf{Relative error} \\
\midrule
0 & 0.05 \\
1 & 0.08 \\
2 & 0.12 \\
3 & 0.15 \\
4 & 0.18 \\
5 & 0.22 \\
\bottomrule
\end{tabular}
\end{table}

Errors grow with depth as feature complexity increases.

\subsection{Energy Preservation Validation}

We verify that the $\sqrt{N/r}$ scaling preserves activation magnitudes:

\begin{table}[H]
\centering
\small
\begin{tabular}{ll}
\toprule
\textbf{Config} & \textbf{Mean activation std} \\
\midrule
Baseline & 1.00 \\
Without scaling & 0.72 \\
With scaling & 1.00 \\
\bottomrule
\end{tabular}
\end{table}

\subsection{Comparison with Fixed Rank Projection}

\textbf{Table 8.3: DCT vs random orthonormal basis.}

\begin{table}[H]
\centering
\small
\begin{tabular}{lll}
\toprule
\textbf{Basis} & \textbf{D acc} & \textbf{R acc} \\
\midrule
Random orthonormal & 87.0 & 74.0 \\
\textbf{DCT} & \textbf{89.91} & \textbf{77.35} \\
\bottomrule
\end{tabular}
\end{table}

DCT's smoothness assumption matches attention better.

\subsection{Failure Analysis}

No systematic failures observed.

\subsection{Summary of Insights}

\begin{enumerate}
\item \textbf{Encoder benefit is real} -- 2$\times$ speedup at $N=128$.
\item \textbf{Decoder fallback necessary} -- no benefit for GPT-2.
\item \textbf{DCT basis optimal} -- beats random orthonormal.
\item \textbf{Energy scaling works} -- preserves activation magnitude.
\item \textbf{Parameter count increases} -- not a compression method.
\end{enumerate}